% ── SECTION 3: THE RELATIONAL COHERENCE THEOREM (~1,200 words) ───────────────
% Formal statement and proof of the four propositions.

\section{The Relational Coherence Theorem}
\label{sec:rct}

\subsection{Core Definitions}

\begin{definition}[Relational Density Matrix]
For composite system $SE$ in pure state $\ket{\psi_{SE}}$, the
\emph{relational density matrix} of $S$ with respect to $E$ is
\begin{equation}
  \rhoSE = \mathrm{Tr}_E\bigl[\ket{\psi_{SE}}\bra{\psi_{SE}}\bigr].
  \label{eq:rho-def}
\end{equation}
\end{definition}

\begin{definition}[Affinity Tensor]
The \emph{Affinity Tensor} $\alpha(S,E)$ is the matrix of off-diagonal
coherence magnitudes of $\rhoSE$:
\begin{equation}
  \alpha(S,E)_{ij} = \bigl|\rho(S|E)_{ij}\bigr|, \quad i \neq j.
  \label{eq:alpha-def}
\end{equation}
The scalar affinity $\bar\alpha = \|\alpha\|_F / \|\alpha_{\max}\|_F$
normalizes to $[0,1]$, where $\bar\alpha = 0$ denotes full decoherence
and $\bar\alpha = 1$ denotes maximal entanglement.
\end{definition}

\subsection{The Four Propositions}

% TODO: state each as a formal Proposition and provide proof sketches.
% Full proofs in Appendix A.

\begin{proposition}[Property Indefiniteness]
No quantum property of system $S$ exists independently of the relational
context established by its interaction with $E$.
\label{prop:indefiniteness}
\end{proposition}

\begin{proof}[Proof sketch]
By Bell's theorem~\cite{Bell1964}, any assignment of definite values to
observables prior to measurement requires hidden variables, which are
excluded by the CHSH experiments~\cite{CHSH1969,Aspect1982}.
See Appendix~\ref{app:rct-proof} for formal derivation.
\end{proof}

\begin{proposition}[Relational Determination]
All definite properties of $S$ emerge through its relational
network~$\Gamma_S = \{(S,E_i) : E_i \text{ interacts with } S\}$.
\label{prop:determination}
\end{proposition}

\begin{proposition}[Coherence--Affinity Correspondence]
The degree of relational participation in being is measured by~$\bar\alpha$:
maximal entanglement ($\bar\alpha = 1$) corresponds to full relational
existence; full decoherence ($\bar\alpha = 0$) corresponds to
classical separability.
\label{prop:correspondence}
\end{proposition}

\begin{proposition}[Conservation of Relational Potential]
In a closed system, $\mathrm{Tr}[\rho^2] = 1$ is conserved under unitary
evolution; total quantum information is neither created nor destroyed.
\label{prop:conservation}
\end{proposition}

\begin{proof}[Proof sketch]
Follows directly from unitarity of time evolution: $\rho(t) = U(t)\rho(0)U^\dagger(t)$
preserves $\mathrm{Tr}[\rho^2]$.
\end{proof}

% TODO: Add figure showing alpha dynamics under a simple decoherence model.
